\section*{Capítulo \ref{ch:b1:non-inertial-frames} --- Sistemas no inerciales; dinámica en la Tierra}

\begin{solution}{exo:b1:non-inertial-frames:1}
$\tan\alpha = a/g = 0.255$, $\alpha = 14^\circ$ (la plomada se queda
atrás); $T = m\sqrt{g^2 + a^2} = 1.03\,mg$. Frenando: el mismo ángulo y la
misma tensión, con la plomada inclinada hacia delante.
\end{solution}

\begin{solution}{exo:b1:non-inertial-frames:2}
$\Omega = 2\pi/4.0 = \qty{1.57}{rad/s}$; $\Omega^2r = \qty{7.4}{m/s^2} = 0.75g$,
dirigida hacia fuera en el sistema giratorio; el asiento (con el
rozamiento y la sujeción) aporta la fuerza real hacia dentro.
\end{solution}

\begin{solution}{exo:b1:non-inertial-frames:3}
$\Omega = 2\pi/86164 = \qty{7.29e-5}{rad/s}$; $\Omega^2R = \qty{0.034}{m/s^2}$,
un $0.35\%$ de $g$. Salir despedido exige $\Omega^2R = g$: $T = 2\pi\sqrt{R/g} =
\qty{84}{min}$ --- diecisiete veces más rápido que ahora.
\end{solution}

\begin{solution}{exo:b1:non-inertial-frames:4}
$a_C = 2\Omega v\sin\lambda = 2 \times 7.29\times10^{-5} \times 27.8 \times 0.707
= \qty{2.9e-3}{m/s^2}$; $F = \qty{1.1}{kN}$. Hacia el norte: empujado al
este, contra el carril derecho (el oriental); hacia el este: empujado al
sur, otra vez contra el carril derecho.
\end{solution}

\begin{solution}{exo:b1:non-inertial-frames:5}
$z = \Omega^2r^2/2g$: \qty{2.0}{mm} a \qty{2}{rad/s} y \qty{5.1}{cm} a
\qty{10}{rad/s}. Un paraboloide concentra un haz paralelo en un punto: el
mercurio giratorio es un espejo parabólico perfecto, barato y que se pule
solo (aunque no puede inclinarse).
\end{solution}

\begin{solution}{exo:b1:non-inertial-frames:6}
La aceleración centrífuga $\Omega^2R\cos\lambda$ apunta hacia fuera del eje;
su componente perpendicular a la vertical local vale $\Omega^2R\cos\lambda
\sin\lambda$, que inclina $\vect g$ en $\delta \approx \Omega^2R\sin\lambda\cos\lambda/g
= 0.034 \times 0.5/9.81 = \qty{1.7e-3}{rad} = 6'$. En el ecuador, la
rotación reduce $g$ en $\Omega^2R = \qty{0.034}{m/s^2}$ (el achatamiento
añade más).
\end{solution}

\begin{solution}{exo:b1:non-inertial-frames:7}
$\ddot x_E = 2\Omega gt\cos\lambda$: $x_E = \tfrac13\Omega g\cos\lambda\,t^3$; con $t =
\sqrt{2h/g} = \qty{4.5}{s}$: $x_E = \tfrac13 \times 7.29\times10^{-5} \times 9.81
\times 0.707 \times 92 = \qty{1.6}{cm}$. Hacia el este: la boca del pozo se
mueve hacia el este más deprisa (radio mayor) que el fondo, y la piedra
conserva ese exceso --- o, en el sistema giratorio, la \omterm{thm:b1:non-inertial-frames:dynamics}{fuerza de Coriolis}
sobre una velocidad descendente apunta al este.
\end{solution}

\begin{solution}{exo:b1:non-inertial-frames:8}
$T_F = 23.93/\sin\lambda$: París \qty{31.8}{h}; polo \qty{23.9}{h}; ecuador
infinito (sin precesión). París: $360^\circ/31.8 = 11.3^\circ$ por hora.
\end{solution}

\begin{solution}{exo:b1:non-inertial-frames:9}
$a_C = 2\Omega v\sin\lambda = \qty{2.1e-3}{m/s^2}$, hacia la derecha. El aire
que converge de todos lados se desvía a la derecha, así que entra en
espiral en sentido antihorario (visto desde arriba): el ciclón del norte.
Un desagüe se vacía en segundos y un tornado se forma en minutos: la
desviación de Coriolis en esos tiempos es despreciable frente a cualquier
remolino inicial.
\end{solution}

\begin{solution}{exo:b1:non-inertial-frames:10}
En el sistema del coche, la gravedad aparente $\vect g - \vect a$ se
inclina hacia atrás. El aire, más pesado, ``cae'' hacia atrás; el globo,
más ligero que el aire que desaloja, es empujado según la vertical
aparente --- se inclina \emph{hacia delante}, al contrario que una plomada.
\end{solution}

\begin{solution}{exo:b1:non-inertial-frames:11}
$\Omega^2R = g$: $\Omega = \qty{0.31}{rad/s}$, $T = \qty{20}{s}$. $a_C = 2\Omega v
= \qty{0.94}{m/s^2}$, alrededor de $0.1g$: caminando a favor de la
rotación se siente un $10\%$ más pesado (la \omterm{thm:b1:non-inertial-frames:dynamics}{fuerza de Coriolis} apunta
hacia fuera) y en contra, más ligero. Una pelota que se deja caer: en el
\omterm{thm:b1:newton-dynamics:laws}{sistema inercial} va recta mientras el suelo se curva hacia arriba; durante
el tiempo de caída $\sqrt{2h/g} = \qty{0.64}{s}$, la desviación es del
orden de $\tfrac13\Omega\,g\,t^3 \cdot 2 \approx 2 \times 0.31 \times
9.81 \times 0.26/3 \approx \qty{0.5}{m}$ --- unos decímetros, en contra de
la rotación.
\end{solution}

\begin{solution}{exo:b1:non-inertial-frames:12}
$GM/(d - r)^2 - GM/d^2 \approx 2GMr/d^3$ (a primer orden en $r/d$). Luna:
$2 \times 6.67\times10^{-11} \times 7.35\times10^{22} \times 6.37\times10^6/
(3.84\times10^8)^3 = \qty{1.1e-6}{m/s^2}$; Sol: $2 \times 1.33\times10^{20}
\times 6.37\times10^6/(1.50\times10^{11})^3 = \qty{5.0e-7}{m/s^2}$; cociente
$0.46$: las mareas solares son casi la mitad de las lunares, aunque la
atracción del Sol sea $180$ veces mayor --- las mareas van como $M/d^3$.
\end{solution}

\begin{solution}{pb:b1:non-inertial-frames:1}
\textbf{1.} $T = 2\pi\sqrt{67/9.81} = \qty{16.4}{s}$; $x_0/\ell = 0.045$
($2.6^\circ$): pequeña.

\textbf{2.} $v_{\max} = x_0\omega_0 = 3.0 \times 0.383 = \qty{1.15}{m/s}$; altura
$x_0^2/2\ell = \qty{6.7}{cm}$.

\textbf{3.} $T = mg + mv_{\max}^2/\ell = mg(1 + 0.002)$: $1.002\,mg$.

\textbf{4.} $\tfrac12\rho C_xSv_{\max}^2 = 0.5 \times 1.2 \times 0.05 \times 1.3 =
\qty{0.04}{N}$, frente a $m\omega_0^2x_0 = \qty{12}{N}$: trescientas veces
menor y, aun así, bastante para detener el péndulo en unas horas.

\textbf{5.} Un cable largo da un periodo largo, y con él una lenteja lenta
y un arrastre relativo pequeño, además de hacer visible la precesión por
oscilación; una lenteja pesada almacena mucha energía para que el
arrastre la coma --- horas de oscilación.

\textbf{6.} El peso (vertical) y la tensión (en el plano del cable y la
vertical, es decir en el plano de oscilación): ninguna fuerza tiene
componente perpendicular a ese plano.

\textbf{7.} El plano es fijo entre las estrellas; el hielo gira debajo
hacia el este con $\Omega$, así que el observador ve girar el plano en
sentido horario, de vuelta al inicio tras un día sidéreo, \qty{23.93}{h}.

\textbf{8.} Cada media oscilación traza una cuerda algo girada: una
roseta; $86164/8.2 \approx 10500$ medias oscilaciones, cada una $0.034^\circ$ más allá.

\textbf{9.} $\vect\Omega = \Omega(0, \cos\lambda, \sin\lambda)$.

\textbf{10.} Con $\vect v' = (\dot x, \dot y, 0)$:
$\vect\Omega\wedge\vect v' = \Omega(-\sin\lambda\,\dot y,\ \sin\lambda\,\dot x,\ -\cos\lambda\,\dot x)$,
de modo que $-2m\vect\Omega\wedge\vect v'$ tiene las componentes horizontales
$2m\Omega\sin\lambda\,(\dot y, -\dot x)$: solo interviene $\Omega\sin\lambda$.

\textbf{11.} $\ddot x = -\omega_0^2x + 2\omega_F\dot y$, $\ddot y = -\omega_0^2y -
2\omega_F\dot x$; $\omega_F = 7.29\times10^{-5} \times 0.753 = \qty{5.5e-5}{rad/s}
\ll \omega_0 = \qty{0.38}{rad/s}$.

\textbf{12.} $\ddot u = -\omega_0^2u + 2\omega_F(\dot y - j\dot x) = -\omega_0^2u -
2j\omega_F\dot u$.

\textbf{13.} $-\alpha^2 - 2\omega_F\alpha + \omega_0^2 = 0$: $\alpha = -\omega_F \pm
\sqrt{\omega_F^2 + \omega_0^2} \approx -\omega_F \pm \omega_0$.

\textbf{14.} El paréntesis es una oscilación plana (una recta fija por el
origen si $A = B$ son reales); el factor $\eu^{-j\omega_Ft}$ gira el conjunto
en sentido horario con $\omega_F$: el plano de oscilación precesa en sentido
horario con $\Omega\sin\lambda$.

\textbf{15.} $T_F = 2\pi/\omega_F = \qty{1.14e5}{s} = \qty{31.8}{h}$; $11.3^\circ$
por hora.

\textbf{16.} $\sin\lambda = 0$: la componente vertical local de $\vect\Omega$
se anula; la \omterm{thm:b1:non-inertial-frames:dynamics}{fuerza de Coriolis} horizontal es nula para el movimiento horizontal.

\textbf{17.} $2m\omega_Fv_{\max} = 2 \times 28 \times 5.5\times10^{-5} \times 1.15 =
\qty{3.5}{mN}$; el peso vale \qty{275}{N} ($10^{-5}$ de él) y la fuerza
recuperadora $m\omega_0^2x_0 = \qty{12}{N}$: tres mil veces menor.

\textbf{18.} $\omega_FT/2 = 5.5\times10^{-5} \times 8.2 = \qty{4.5e-4}{rad}$:
$x_0 \times 4.5\times10^{-4} = \qty{1.4}{mm}$ por media oscilación.

\textbf{19.} Con las clavijas a un centímetro unas de otras en el anillo,
cae una cada $\sim 7$ medias oscilaciones, alrededor de un minuto.

\textbf{20.} Un empujón con la mano da una velocidad lateral: la lenteja
describe entonces una elipse, cuya propia precesión lenta (efecto de la
\omterm{def:b1:wave-propagation:signal}{amplitud} finita) enmascara la de Foucault.

\textbf{21.} Un impulso radial añade energía a lo largo de la oscilación
sin añadir velocidad transversal, así que el plano queda intacto; una
componente lateral llevaría el péndulo a una elipse y falsearía o
escondería la precesión.

\textbf{22.} El ritmo de precesión mide $\Omega\sin\lambda$, la componente
vertical de la rotación terrestre; conocida $\Omega$ por la astronomía, el
ritmo da la latitud (y a la inversa).

\textbf{23.} $\sin\lambda = 23.93/48 = 0.499$: $\lambda = 30^\circ$.

\textbf{24.} Su \omterm{def:b1:angular-momentum:L}{momento angular} se conserva, de modo que su eje queda fijo
entre las estrellas mientras la Tierra gira debajo: el girocompás, que
encuentra el norte verdadero sin imán alguno.

\textbf{25.} La \omterm{thm:b1:non-inertial-frames:dynamics}{fuerza de Coriolis} $-2m\vect\Omega\wedge\vect v'$; su parte
horizontal crece como $\sin\lambda$; una vuelta en $23.93\,\text{h}/\sin\lambda$,
\qty{31.8}{h} en París.
\end{solution}
